\subsection{Compile-Time Scope Disambiguation}

Python name lookup feels dynamic because values are resolved while the program runs. But the classification of many names is decided earlier, when Python compiles the function body.

Before a function runs, Python has already inspected its syntax and decided which names are local, which names are global-style lookups, which names are free variables from enclosing scopes, and which names must be stored in closure cells.

This distinction is essential:

\begin{itemize}
    \item name \textbf{classification} is mostly static and syntax-driven;
    \item object \textbf{value lookup} is dynamic and happens at runtime.
\end{itemize}

That is why a function can fail with \texttt{UnboundLocalError} before it ever considers a global name with the same spelling. The compiler has already decided that the name is local.

\subsubsection{How the Python Compiler Scans Syntax Trees for Assignment Targets}

Python first parses source code into a syntax tree. That tree contains assignment targets, reads, function definitions, calls, literals, and other syntactic structures.

The standard \texttt{ast} module lets us inspect this representation.

\begin{verbatim}
import ast

source = """
answer = 42

def broken():
    print(answer)
    answer = 100
"""

tree = ast.parse(source)

print(f"type(tree): {type(tree)}")
print()
print(ast.dump(tree, indent=4))
\end{verbatim}

The output is long, but the important part contains two different uses of \texttt{answer} inside \texttt{broken()}:

\begin{verbatim}
Name(id='answer', ctx=Load())
Name(id='answer', ctx=Store())
\end{verbatim}

The \texttt{Load()} context means the name is being read. The \texttt{Store()} context means the name is an assignment target.

The AST object can be inspected.

\begin{verbatim}
import ast

tree = ast.parse("x = 42")

print(f"type(tree): {type(tree)}")
print()

selected_names = [
    "body",
    "type_ignores",
    "_fields",
    "__class__",
]

for name in selected_names:
    print(f"{name:<14} {name in dir(tree)}")

print()
print(f"tree.body: {tree.body}")
\end{verbatim}

Possible output:

\begin{verbatim}
type(tree): <class 'ast.Module'>

body           True
type_ignores   True
_fields        True
__class__      True

tree.body: [<ast.Assign object at 0x...>]
\end{verbatim}

The first statement is an \texttt{ast.Assign} object.

\begin{verbatim}
import ast

tree = ast.parse("x = 42")
assign_node = tree.body[0]

print(f"assign_node:       {assign_node}")
print(f"type(assign_node): {type(assign_node)}")
print()

selected_names = [
    "targets",
    "value",
    "lineno",
    "col_offset",
    "_fields",
]

for name in selected_names:
    print(f"{name:<14} {name in dir(assign_node)}")

print()
print(f"targets: {assign_node.targets}")
print(f"value:   {assign_node.value}")
\end{verbatim}

Possible output:

\begin{verbatim}
assign_node:       <ast.Assign object at 0x...>
type(assign_node): <class 'ast.Assign'>

targets        True
value          True
lineno         True
col_offset     True
_fields        True

targets: [<ast.Name object at 0x...>]
value:   <ast.Constant object at 0x...>
\end{verbatim}

The assignment target is an \texttt{ast.Name} object with \texttt{Store} context.

\begin{verbatim}
import ast

tree = ast.parse("x = 42")
target_node = tree.body[0].targets[0]

print(f"type(target_node): {type(target_node)}")
print()

selected_names = [
    "id",
    "ctx",
    "lineno",
    "col_offset",
    "_fields",
]

for name in selected_names:
    print(f"{name:<12} {name in dir(target_node)}")

print()
print(f"target_node.id:  {target_node.id!r}")
print(f"target_node.ctx: {target_node.ctx}")
print(f"type(ctx):       {type(target_node.ctx)}")
\end{verbatim}

Possible output:

\begin{verbatim}
type(target_node): <class 'ast.Name'>

id           True
ctx          True
lineno       True
col_offset   True
_fields      True

target_node.id:  'x'
target_node.ctx: <ast.Store object at 0x...>
type(ctx):       <class 'ast.Store'>
\end{verbatim}

Now compare a name read.

\begin{verbatim}
import ast

tree = ast.parse("print(x)")
call_node = tree.body[0].value
name_node = call_node.args[0]

print(f"type(name_node): {type(name_node)}")
print(f"name_node.id:    {name_node.id!r}")
print(f"name_node.ctx:   {name_node.ctx}")
print(f"type(ctx):       {type(name_node.ctx)}")
\end{verbatim}

Possible output:

\begin{verbatim}
type(name_node): <class 'ast.Name'>
name_node.id:    'x'
name_node.ctx:   <ast.Load object at 0x...>
type(ctx):       <class 'ast.Load'>
\end{verbatim}

The compiler uses this kind of syntactic information to decide scope categories. A name that appears as an assignment target inside a function is local by default, unless a \texttt{global} or \texttt{nonlocal} declaration changes that rule.

This is true even if the assignment is textually after the read.

\begin{verbatim}
answer = 42

def broken():
    print(answer)
    answer = 100

try:
    broken()
except UnboundLocalError as err:
    print(f"error type: {type(err)}")
    print(f"message:    {err}")
\end{verbatim}

Possible output:

\begin{verbatim}
error type: <class 'UnboundLocalError'>
message:    cannot access local variable 'answer' where it is not associated with a value
\end{verbatim}

The compiler saw the assignment target \texttt{answer} in the function body. Therefore, \texttt{answer} became a local variable for that function. The earlier read tries to read the local slot before a value has been stored there.

A dead-looking assignment still affects classification.

\begin{verbatim}
answer = 42

def broken():
    if False:
        answer = 100

    return answer

try:
    print(broken())
except UnboundLocalError as err:
    print(f"error type: {type(err)}")
    print(f"message:    {err}")
\end{verbatim}

Possible output:

\begin{verbatim}
error type: <class 'UnboundLocalError'>
message:    cannot access local variable 'answer' where it is not associated with a value
\end{verbatim}

The branch does not execute, but the assignment target exists in the function syntax. The compiler does not need the branch to run before it classifies \texttt{answer} as local.

The code object reveals the decision.

\begin{verbatim}
answer = 42

def broken():
    if False:
        answer = 100

    return answer

print(f"co_varnames: {broken.__code__.co_varnames}")
print(f"co_names:    {broken.__code__.co_names}")
\end{verbatim}

Possible output:

\begin{verbatim}
co_varnames: ('answer',)
co_names:    ()
\end{verbatim}

The name \texttt{answer} is a local variable of \texttt{broken()}, even though the assignment never runs.

The rule is:

\begin{quote}
For ordinary function bodies, assignment syntax classifies a name as local for the whole function unless \texttt{global} or \texttt{nonlocal} overrides that classification.
\end{quote}

\subsubsection{Pre-Determining Local Scope Allocation Flags via Symbol Tables Prior to Execution}

The standard \texttt{symtable} module exposes Python's symbol-table analysis. This is closer to the compiler's scope reasoning than plain AST inspection.

\begin{verbatim}
import symtable

source = """
answer = 42

def broken():
    print(answer)
    answer = 100
"""

table = symtable.symtable(source, "demo.py", "exec")

print(f"table:       {table}")
print(f"type(table): {type(table)}")
\end{verbatim}

Possible output:

\begin{verbatim}
table:       <SymbolTable for module demo.py>
type(table): <class 'symtable.SymbolTable'>
\end{verbatim}

The symbol table object can be inspected.

\begin{verbatim}
import symtable

source = """
answer = 42

def broken():
    print(answer)
    answer = 100
"""

table = symtable.symtable(source, "demo.py", "exec")

print(f"type(table): {type(table)}")
print()

selected_names = [
    "get_name",
    "get_type",
    "get_symbols",
    "get_children",
    "lookup",
    "is_nested",
]

for name in selected_names:
    print(f"{name:<16} {name in dir(table)}")

print()
print(f"table name: {table.get_name()}")
print(f"table type: {table.get_type()}")
\end{verbatim}

Possible output:

\begin{verbatim}
type(table): <class 'symtable.SymbolTable'>

get_name         True
get_type         True
get_symbols      True
get_children     True
lookup           True
is_nested        True

table name: top
table type: module
\end{verbatim}

The module-level symbol table has a child table for the function.

\begin{verbatim}
import symtable

source = """
answer = 42

def broken():
    print(answer)
    answer = 100
"""

module_table = symtable.symtable(source, "demo.py", "exec")
children = module_table.get_children()

print(f"children:       {children}")
print(f"type(children): {type(children)}")
print()

func_table = children[0]

print(f"func name: {func_table.get_name()}")
print(f"func type: {func_table.get_type()}")
\end{verbatim}

Possible output:

\begin{verbatim}
children:       [<Function SymbolTable for broken in demo.py>]
type(children): <class 'list'>

func name: broken
func type: function
\end{verbatim}

The child list is a normal list.

\begin{verbatim}
import symtable

source = """
def broken():
    answer = 100
"""

module_table = symtable.symtable(source, "demo.py", "exec")
children = module_table.get_children()

print(f"type(children): {type(children)}")
print()

selected_names = [
    "append",
    "extend",
    "pop",
    "__len__",
    "__getitem__",
    "__iter__",
]

for name in selected_names:
    print(f"{name:<14} {name in dir(children)}")

print()
print(f"len(children): {len(children)}")
\end{verbatim}

Possible output:

\begin{verbatim}
type(children): <class 'list'>

append         True
extend         True
pop            True
__len__        True
__getitem__    True
__iter__       True

len(children): 1
\end{verbatim}

Now inspect the symbol for \texttt{answer} inside the function.

\begin{verbatim}
import symtable

source = """
answer = 42

def broken():
    print(answer)
    answer = 100
"""

module_table = symtable.symtable(source, "demo.py", "exec")
func_table = module_table.get_children()[0]

symbol = func_table.lookup("answer")

print(f"symbol:       {symbol}")
print(f"type(symbol): {type(symbol)}")
print()

selected_names = [
    "get_name",
    "is_local",
    "is_global",
    "is_referenced",
    "is_assigned",
    "is_free",
    "is_parameter",
]

for name in selected_names:
    print(f"{name:<18} {name in dir(symbol)}")

print()
print(f"name:          {symbol.get_name()}")
print(f"is_local:      {symbol.is_local()}")
print(f"is_global:     {symbol.is_global()}")
print(f"is_referenced: {symbol.is_referenced()}")
print(f"is_assigned:   {symbol.is_assigned()}")
print(f"is_free:       {symbol.is_free()}")
\end{verbatim}

Possible output:

\begin{verbatim}
symbol:       <symbol 'answer'>
type(symbol): <class 'symtable.Symbol'>

get_name           True
is_local           True
is_global          True
is_referenced      True
is_assigned        True
is_free            True
is_parameter       True

name:          answer
is_local:      True
is_global:     False
is_referenced: True
is_assigned:   True
is_free:       False
\end{verbatim}

The symbol table reports that \texttt{answer} is local and assigned inside the function. This happens before the function is ever called.

Compare a function that only reads \texttt{answer}.

\begin{verbatim}
import symtable

source = """
answer = 42

def show():
    print(answer)
"""

module_table = symtable.symtable(source, "demo.py", "exec")
func_table = module_table.get_children()[0]
symbol = func_table.lookup("answer")

print(f"name:          {symbol.get_name()}")
print(f"is_local:      {symbol.is_local()}")
print(f"is_global:     {symbol.is_global()}")
print(f"is_referenced: {symbol.is_referenced()}")
print(f"is_assigned:   {symbol.is_assigned()}")
print(f"is_free:       {symbol.is_free()}")
\end{verbatim}

Possible output:

\begin{verbatim}
name:          answer
is_local:      False
is_global:     True
is_referenced: True
is_assigned:   False
is_free:       False
\end{verbatim}

Now \texttt{answer} is classified as global from the function's point of view.

The compiled code objects agree with the symbol table.

\begin{verbatim}
answer = 42

def show():
    print(answer)

def broken():
    print(answer)
    answer = 100

print("show")
print(f"co_varnames: {show.__code__.co_varnames}")
print(f"co_names:    {show.__code__.co_names}")
print()

print("broken")
print(f"co_varnames: {broken.__code__.co_varnames}")
print(f"co_names:    {broken.__code__.co_names}")
\end{verbatim}

Possible output:

\begin{verbatim}
show
co_varnames: ()
co_names:    ('print', 'answer')

broken
co_varnames: ('answer',)
co_names:    ('print',)
\end{verbatim}

The symbol-table decision affects the bytecode instruction choice.

\begin{verbatim}
import dis

answer = 42

def show():
    print(answer)

def broken():
    print(answer)
    answer = 100

print("show:")
dis.dis(show)

print()
print("broken:")
dis.dis(broken)
\end{verbatim}

Possible output excerpt:

\begin{verbatim}
show:
... LOAD_GLOBAL              ... (answer)

broken:
... LOAD_FAST                ... (answer)
... STORE_FAST               ... (answer)
\end{verbatim}

The exact bytecode format depends on Python version, but the difference remains:

\begin{itemize}
    \item global-style names use global lookup instructions;
    \item local names use local-slot instructions.
\end{itemize}

The symbol table is therefore a compile-time planning artifact. It pre-determines how names will be accessed when the function eventually runs.

\subsubsection{Static Name Binding Invariants vs. Dynamic Late-Binding Value Resolution}

The classification of a name is static for a compiled function. The value found through that classification may still be dynamic.

A function that reads a global name does not freeze that global value when the function is defined.

\begin{verbatim}
message = "D'oh!"

def show_message():
    print(message)

show_message()

message = "No soup for you!"

show_message()
\end{verbatim}

Possible output:

\begin{verbatim}
D'oh!
No soup for you!
\end{verbatim}

The function \texttt{show\_message()} performs a global lookup each time it is called. It does not store the original string as a hidden constant.

The code object shows that \texttt{message} is a global-style name.

\begin{verbatim}
message = "D'oh!"

def show_message():
    print(message)

print(f"co_varnames: {show_message.__code__.co_varnames}")
print(f"co_names:    {show_message.__code__.co_names}")
\end{verbatim}

Possible output:

\begin{verbatim}
co_varnames: ()
co_names:    ('print', 'message')
\end{verbatim}

The classification is fixed:

\begin{verbatim}
message is looked up as a global name
\end{verbatim}

The value is not fixed:

\begin{verbatim}
whatever object the module namespace binds to message at call time is used
\end{verbatim}

This is late binding of values.

The same idea appears with functions.

\begin{verbatim}
def action():
    return "first action"

def run_action():
    print(action())

run_action()

def action():
    return "second action"

run_action()
\end{verbatim}

Possible output:

\begin{verbatim}
first action
second action
\end{verbatim}

The function \texttt{run\_action()} looks up the global name \texttt{action} when it runs. After the second \texttt{def action():} executes, the global name \texttt{action} refers to a different function object.

The function object can be inspected before and after rebinding.

\begin{verbatim}
def action():
    return "first action"

old_action = action

def action():
    return "second action"

print(f"old_action: {old_action}")
print(f"action:     {action}")
print(f"same object: {old_action is action}")
print()
print(f"old_action(): {old_action()}")
print(f"action():     {action()}")
\end{verbatim}

Possible output:

\begin{verbatim}
old_action: <function action at 0x...>
action:     <function action at 0x...>
same object: False

old_action(): first action
action():     second action
\end{verbatim}

The global name was rebound. Code that looks up the global name later sees the new object.

This contrasts with local classification.

\begin{verbatim}
message = "global message"

def broken():
    print(message)
    message = "local message"

try:
    broken()
except UnboundLocalError as err:
    print(f"error type: {type(err)}")
    print(f"message:    {err}")
\end{verbatim}

Possible output:

\begin{verbatim}
error type: <class 'UnboundLocalError'>
message:    cannot access local variable 'message' where it is not associated with a value
\end{verbatim}

Even though the global \texttt{message} exists at runtime, the compiled function does not use the global lookup path for \texttt{message}. The local classification is fixed by syntax.

The invariant is:

\begin{quote}
The compiler decides the namespace access category. Runtime execution retrieves the current object from that category.
\end{quote}

A default argument behaves differently from a global lookup. The default object is stored when the function is defined.

\begin{verbatim}
message = "D'oh!"

def show_default(text=message):
    print(text)

show_default()

message = "No soup for you!"

show_default()
show_default(message)
\end{verbatim}

Possible output:

\begin{verbatim}
D'oh!
D'oh!
No soup for you!
\end{verbatim}

The default value captured the old object at definition time. The global lookup in a function body would see the new object at call time.

Compare both in one example.

\begin{verbatim}
message = "D'oh!"

def show_both(text=message):
    print(f"default text: {text}")
    print(f"global name:  {message}")

show_both()

message = "No soup for you!"

print()
show_both()
\end{verbatim}

Possible output:

\begin{verbatim}
default text: D'oh!
global name:  D'oh!

default text: D'oh!
global name:  No soup for you!
\end{verbatim}

The default \texttt{text} is stored in \texttt{\_\_defaults\_\_}. The name \texttt{message} inside the function body is looked up dynamically as a global.

\begin{verbatim}
message = "D'oh!"

def show_both(text=message):
    print(text)
    print(message)

print(f"defaults:    {show_both.__defaults__}")
print(f"co_varnames: {show_both.__code__.co_varnames}")
print(f"co_names:    {show_both.__code__.co_names}")
\end{verbatim}

Possible output:

\begin{verbatim}
defaults:    ("D'oh!",)
co_varnames: ('text',)
co_names:    ('print', 'message')
\end{verbatim}

This example separates three ideas:

\begin{itemize}
    \item \texttt{text} is a local parameter;
    \item its default object was stored at definition time;
    \item \texttt{message} in the body is a global lookup resolved at call time.
\end{itemize}

So Python is not simply ``dynamic'' or ``static'' in one flat sense. Its behavior is split:

\begin{center}
\begin{tabular}{|p{0.38\linewidth}|p{0.46\linewidth}|}
\hline
\textbf{Mechanism} & \textbf{When it is decided} \\
\hline
Local/global/free/built-in classification & compile time / function creation path \\
\hline
Default argument object creation & definition time, when \texttt{def} executes \\
\hline
Global value lookup inside function body & call time, when instruction executes \\
\hline
Local value availability & call time, depending on executed assignments \\
\hline
Closure cell value lookup & call time, through preserved cell object \\
\hline
\end{tabular}
\end{center}

The important practical rule is:

\begin{quote}
The place Python will search is decided early. The object found in that place may still change before the search happens.
\end{quote}

\subsubsection{The Loop Variable Closure Trap: Why Nested Functions May See the Final Loop Value}

A common closure trap appears when nested functions are created inside a loop.

\begin{verbatim}
funcs = []

for i in range(3):
    def show_i():
        return i

    funcs.append(show_i)

for fn in funcs:
    print(fn())
\end{verbatim}

Possible output:

\begin{verbatim}
2
2
2
\end{verbatim}

Many beginners expect:

\begin{verbatim}
0
1
2
\end{verbatim}

But all three functions refer to the same name \texttt{i}. They do not each receive a frozen copy of the loop value.

At module level, \texttt{i} is a global name. After the loop finishes, \texttt{i} is bound to \texttt{2}. Each function performs a global lookup of \texttt{i} when called.

This can be inspected.

\begin{verbatim}
funcs = []

for i in range(3):
    def show_i():
        return i

    funcs.append(show_i)

fn = funcs[0]

print(f"fn(): {fn()}")
print(f"i after loop: {i}")
print()
print(f"fn.__code__.co_varnames: {fn.__code__.co_varnames}")
print(f"fn.__code__.co_names:    {fn.__code__.co_names}")
print(f"fn.__closure__:          {fn.__closure__}")
\end{verbatim}

Possible output:

\begin{verbatim}
fn(): 2
i after loop: 2

fn.__code__.co_varnames: ()
fn.__code__.co_names:    ('i',)
fn.__closure__:          None
\end{verbatim}

At module level, the nested-looking function is not inside another function scope. It uses global lookup for \texttt{i}.

Inside an outer function, the trap uses closure cells.

\begin{verbatim}
def make_funcs():
    funcs = []

    for i in range(3):
        def show_i():
            return i

        funcs.append(show_i)

    return funcs

funcs = make_funcs()

for fn in funcs:
    print(fn())
\end{verbatim}

Possible output:

\begin{verbatim}
2
2
2
\end{verbatim}

Now \texttt{i} is not global. It is a local variable of \texttt{make\_funcs()} that becomes a closure cell because \texttt{show\_i()} refers to it.

All three inner functions share that same cell.

\begin{verbatim}
def make_funcs():
    funcs = []

    for i in range(3):
        def show_i():
            return i

        funcs.append(show_i)

    return funcs

funcs = make_funcs()

for index, fn in enumerate(funcs):
    cell = fn.__closure__[0]

    print(f"function {index}")
    print(f"  result:          {fn()}")
    print(f"  closure:         {fn.__closure__}")
    print(f"  freevars:        {fn.__code__.co_freevars}")
    print(f"  cell contents:   {cell.cell_contents}")
    print(f"  cell id:         {id(cell)}")
    print()
\end{verbatim}

Possible output:

\begin{verbatim}
function 0
  result:          2
  closure:         (<cell at 0x...: int object at 0x...>,)
  freevars:        ('i',)
  cell contents:   2
  cell id:         2293100811520

function 1
  result:          2
  closure:         (<cell at 0x...: int object at 0x...>,)
  freevars:        ('i',)
  cell contents:   2
  cell id:         2293100811520

function 2
  result:          2
  closure:         (<cell at 0x...: int object at 0x...>,)
  freevars:        ('i',)
  cell contents:   2
  cell id:         2293100811520
\end{verbatim}

The same cell id appears for all three functions. They share one closure cell for \texttt{i}. After the loop ends, that cell contains the final loop value.

The cell object is a normal inspectable object.

\begin{verbatim}
def make_funcs():
    funcs = []

    for i in range(3):
        def show_i():
            return i

        funcs.append(show_i)

    return funcs

cell = make_funcs()[0].__closure__[0]

print(f"type(cell): {type(cell)}")
print()

selected_names = [
    "cell_contents",
    "__class__",
    "__eq__",
    "__repr__",
]

for name in selected_names:
    print(f"{name:<16} {name in dir(cell)}")

print()
print(f"cell.cell_contents: {cell.cell_contents}")
\end{verbatim}

Possible output:

\begin{verbatim}
type(cell): <class 'cell'>

cell_contents    True
__class__        True
__eq__           True
__repr__         True

cell.cell_contents: 2
\end{verbatim}

The problem is not that the loop calls the functions too late in a vague sense. The exact mechanism is:

\begin{verbatim}
1. The loop repeatedly rebinds the same name i.
2. Each inner function refers to that same variable cell.
3. The functions are called after the loop has finished.
4. The cell now contains the final value, 2.
5. Every function reads that same cell.
\end{verbatim}

A common fix is to capture the current value as a default argument.

\begin{verbatim}
def make_funcs():
    funcs = []

    for i in range(3):
        def show_i(i=i):
            return i

        funcs.append(show_i)

    return funcs

funcs = make_funcs()

for fn in funcs:
    print(fn())
\end{verbatim}

Possible output:

\begin{verbatim}
0
1
2
\end{verbatim}

This works because default argument objects are evaluated when the \texttt{def} statement executes. Each loop iteration executes a new \texttt{def} statement and stores the current object referenced by \texttt{i} as that function's default.

Inspect the functions.

\begin{verbatim}
def make_funcs():
    funcs = []

    for i in range(3):
        def show_i(i=i):
            return i

        funcs.append(show_i)

    return funcs

funcs = make_funcs()

for index, fn in enumerate(funcs):
    print(f"function {index}")
    print(f"  result:       {fn()}")
    print(f"  defaults:     {fn.__defaults__}")
    print(f"  varnames:     {fn.__code__.co_varnames}")
    print(f"  freevars:     {fn.__code__.co_freevars}")
    print(f"  closure:      {fn.__closure__}")
    print()
\end{verbatim}

Possible output:

\begin{verbatim}
function 0
  result:       0
  defaults:     (0,)
  varnames:     ('i',)
  freevars:     ()
  closure:      None

function 1
  result:       1
  defaults:     (1,)
  varnames:     ('i',)
  freevars:     ()
  closure:      None

function 2
  result:       2
  defaults:     (2,)
  varnames:     ('i',)
  freevars:     ()
  closure:      None
\end{verbatim}

Now each function has its own local parameter \texttt{i} with its own stored default. There is no shared closure cell for \texttt{i}.

Another fix is to use a factory function. The factory creates a new enclosing scope for each value.

\begin{verbatim}
def make_one_func(value):
    def show_value():
        return value

    return show_value

funcs = []

for i in range(3):
    funcs.append(make_one_func(i))

for fn in funcs:
    print(fn())
\end{verbatim}

Possible output:

\begin{verbatim}
0
1
2
\end{verbatim}

Each call to \texttt{make\_one\_func(i)} creates a separate frame and a separate closure cell for \texttt{value}.

\begin{verbatim}
def make_one_func(value):
    def show_value():
        return value

    return show_value

funcs = []

for i in range(3):
    funcs.append(make_one_func(i))

for index, fn in enumerate(funcs):
    cell = fn.__closure__[0]

    print(f"function {index}")
    print(f"  result:        {fn()}")
    print(f"  freevars:      {fn.__code__.co_freevars}")
    print(f"  cell contents: {cell.cell_contents}")
    print(f"  cell id:       {id(cell)}")
    print()
\end{verbatim}

Possible output:

\begin{verbatim}
function 0
  result:        0
  freevars:      ('value',)
  cell contents: 0
  cell id:       2293100811520

function 1
  result:        1
  freevars:      ('value',)
  cell contents: 1
  cell id:       2293100811392

function 2
  result:        2
  freevars:      ('value',)
  cell contents: 2
  cell id:       2293100811648
\end{verbatim}

Now the cell ids are different. Each returned function closes over a different cell.

The loop variable closure trap is therefore not a special loop bug. It is a direct consequence of two ordinary rules:

\begin{itemize}
    \item closures preserve variables through cells, not frozen textual values;
    \item the value inside a cell is read when the inner function runs, not when the inner function is created.
\end{itemize}

The compact rule is:

\begin{quote}
If a nested function refers to a loop variable, it usually refers to the variable itself, not a snapshot of its value. Use a default argument or a factory function when each generated function must remember a separate value.
\end{quote}